%============================================================
%  Bd. 16 - Monoide
%============================================================

\documentclass[main.tex]{subfiles}

\ifSubfilesClassLoaded{
  \BandSetupStandalone{B16}
}{
}

\title{Bd. 16 - Monoide}
\author{Martin Kunze}
\date{}
\setcounter{file}{16}

\begin{document}

\title{Bd. 16 - Monoide}
\author{Martin Kunze}
\date{}
\hypersetup{pageanchor=false}
\maketitle
\hypersetup{pageanchor=true}
\setcounter{tocdepth}{3}
\tableofcontents
\setcounter{file}{16}
\renewcommand{\FormulaBandID}{16}

\chapter{Einführung}

Dieser Band entwickelt Monoide als Halbgruppen mit einem neutralen Element.
Darauf aufbauend werden Einheiten, Teilbarkeit, Irreduzibilität,
Primelemente und Morphismen von Monoiden untersucht.

\chapter{Monoide}

\section{Definition des Monoids}

\FormulaDefDeltaK[Begriff des Monoids]{%
  \MonoidAlg{A}{\star}{e}%
}{MonoidDef}{%
  \DeltaRow{Mengen}{A\dsep e\dsep x}
  \DeltaRow{\textbf{Axiome}}{}
  \DeltaPrem{Halbgruppe}{\SemigroupAlg{A}{\star}}
  \DeltaRow{Neutrales Element im Träger}{e\in A}
  \DeltaRow{Linksneutralität}{x\in A\vdash e\star x=x}
  \DeltaRow{Rechtsneutralität}{x\in A\vdash x\star e=x}
  \DeltaRow{\textbf{Neues Symbol}}{}
  \DeltaPrem{Monoid}{\MonoidAlg{A}{\star}{e}}
}

\FormulaAxiomDeltaK[Zugriff auf die Halbgruppe eines Monoids]{%
  \MonoidAlg{A}{\star}{e}\vdash\SemigroupAlg{A}{\star}%
}{MonoidBaseSemigroupAxiom}{%
  \DeltaRow{Mengen}{A\dsep e}
}

\FormulaAxiomDeltaK[Zugriff auf das neutrale Element]{%
  \MonoidAlg{A}{\star}{e}\vdash e\in A%
}{MonoidIdentityCarrierAxiom}{%
  \DeltaRow{Mengen}{A\dsep e}
}

\FormulaAxiomDeltaK[Zugriff auf die Linksneutralität]{%
  \MonoidAlg{A}{\star}{e}\dsep x\in A
  \vdash e\star x=x%
}{MonoidLeftIdentityAxiom}{%
  \DeltaRow{Mengen}{A\dsep e\dsep x}
}

\FormulaAxiomDeltaK[Zugriff auf die Rechtsneutralität]{%
  \MonoidAlg{A}{\star}{e}\dsep x\in A
  \vdash x\star e=x%
}{MonoidRightIdentityAxiom}{%
  \DeltaRow{Mengen}{A\dsep e\dsep x}
}

\section{Das neutrale Element}

\FormulaThmDeltaK[Ein neutraler Kandidat stimmt mit dem neutralen Element überein]{%
  \begin{aligned}[t]
    &\MonoidAlg{A}{\star}{e}\dsep e'\in A\dsep
      \forall x\in A\,(e'\star x=x)\dsep
      \forall x\in A\,(x\star e'=x)\vdash{}\\[-2pt]
    &e=e'
  \end{aligned}
}{MonoidIdentityCandidateEqualsIdentity}{%
  \DeltaRow{Mengen}{A\dsep e\dsep e'\dsep x}
}
\begin{tabproofwide}
    \proofstepwidestar[1]{\MonoidAlg{A}{\star}{e}}{\rA}
    \proofstepwidestar[2]{e'\in A}{\rA}
    \proofstepwidestar[3]{\forall x\in A\,(e'\star x=x)}{\rA}
    \proofstepwidestar[4]{\forall x\in A\,(x\star e'=x)}{\rA}
    \proofstepwidestar[1]{e\in A}{%
      \FormulaRefAuto{MonoidIdentityCarrierAxiom}[ax]{1}}
    \proofstepwidestar[1,3]{e'\star e=e}{\rRE{5,3}}
    \proofstepwide[1,3]{e}{=}{e'\star e}{%
      \FormulaRefAuto{a=b\vdash b=a}{6}}
    \proofstepwide[1,2]{}{=}{e'}{%
      \FormulaRefAuto{MonoidRightIdentityAxiom}[ax]{1,2}}
    \proofstepwidestar[1,2,3]{e=e'}{\rChain{7,8}}
\end{tabproofwide}

\FormulaThmDeltaK[Existenz und Eindeutigkeit des neutralen Elements]{%
  \begin{aligned}[t]
    &\MonoidAlg{A}{\star}{e}\vdash{}\\[-2pt]
    &\exists! e'\in A\,\Bigl(
      \forall x\in A\,(e'\star x=x)\land
      \forall x\in A\,(x\star e'=x)
      \Bigr)
  \end{aligned}
}{MonoidIdentityUnique}{%
  \DeltaRow{Mengen}{A\dsep e\dsep e'\dsep x}
}
\begin{tabproofwide}
    \proofstepwidestar[1]{\MonoidAlg{A}{\star}{e}}{\rA}
    \proofstepwidestar[1]{e\in A}{%
      \FormulaRefAuto{MonoidIdentityCarrierAxiom}[ax]{1}}
    \proofstepwidestar[3]{x\in A}{\rA}
    \proofstepwidestar[1,3]{e\star x=x}{%
      \FormulaRefAuto{MonoidLeftIdentityAxiom}[ax]{1,3}}
    \proofstepwidestar[1,3]{x\star e=x}{%
      \FormulaRefAuto{MonoidRightIdentityAxiom}[ax]{1,3}}
    \proofstepwidestar[1]{\forall x\in A\,(e\star x=x)}{%
      \rUI{\rRI{3,4}}}
    \proofstepwidestar[1]{\forall x\in A\,(x\star e=x)}{%
      \rUI{\rRI{3,5}}}
    \proofstepwidestar[1]{%
      \begin{aligned}[t]
        &e\in A\land\Bigl(
          \forall x\in A\,(e\star x=x)\land{}\\[-2pt]
        &\qquad \forall x\in A\,(x\star e=x)
          \Bigr)
      \end{aligned}}{%
      \rAI{2,\rAI{6,7}}}
    \proofstepwidestar[9]{%
      \begin{aligned}[t]
        &e'\in A\land\Bigl(
          \forall x\in A\,(e'\star x=x)\land{}\\[-2pt]
        &\qquad \forall x\in A\,(x\star e'=x)
          \Bigr)
      \end{aligned}}{\rA}
    \proofstepwidestar[9]{e'\in A}{\rAEa{9}}
    \proofstepwidestar[9]{\forall x\in A\,(e'\star x=x)}{%
      \rAEa{\rAEb{9}}}
    \proofstepwidestar[9]{\forall x\in A\,(x\star e'=x)}{%
      \rAEb{\rAEb{9}}}
    \proofstepwidestar[1,9]{e=e'}{%
      \FormulaRefAuto{MonoidIdentityCandidateEqualsIdentity}[thm]{%
        1,10,11,12}}
    \proofstepwidestar[1]{%
      \begin{aligned}[t]
        &\exists! e'\in A\,\Bigl(
          \forall x\in A\,(e'\star x=x)\land{}\\[-2pt]
        &\qquad \forall x\in A\,(x\star e'=x)
          \Bigr)
      \end{aligned}}{\UEI{8,9,13}}
\end{tabproofwide}

\section{Spezielle Definitionen}

\subsection{Endliche Monoide}

\FormulaDefDeltaK[Begriff des endlichen Monoids]{%
  \FiniteMonoidAlg{A}{\star}{e}%
}{FiniteMonoidDef}{%
  \DeltaRow{Mengen}{A\dsep e}
  \DeltaRow{\textbf{Axiome}}{}
  \DeltaPrem{Monoid}{\MonoidAlg{A}{\star}{e}}
  \DeltaPrem{Endlichkeit}{\Finite{A}}
  \DeltaRow{\textbf{Neues Symbol}}{}
  \DeltaPrem{Endliches Monoid}{\FiniteMonoidAlg{A}{\star}{e}}
}

\FormulaAxiomDeltaK[Zugriff auf die Monoidstruktur]{%
  \FiniteMonoidAlg{A}{\star}{e}\vdash\MonoidAlg{A}{\star}{e}%
}{FiniteMonoidBaseAxiom}{%
  \DeltaRow{Mengen}{A\dsep e}
}

\FormulaAxiomDeltaK[Zugriff auf die Endlichkeit]{%
  \FiniteMonoidAlg{A}{\star}{e}\vdash\Finite{A}%
}{FiniteMonoidFinitenessAxiom}{%
  \DeltaRow{Mengen}{A\dsep e}
}

\subsection{Kommutative Monoide}

\FormulaDefDeltaK[Begriff des kommutativen Monoids]{%
  \CommutativeMonoidAlg{A}{\star}{e}%
}{CommutativeMonoidDef}{%
  \DeltaRow{Mengen}{A\dsep e\dsep x\dsep y}
  \DeltaRow{\textbf{Axiome}}{}
  \DeltaPrem{Monoid}{\MonoidAlg{A}{\star}{e}}
  \DeltaRow{Kommutativität}{x,y\in A\vdash x\star y=y\star x}
  \DeltaRow{\textbf{Neues Symbol}}{}
  \DeltaPrem{Kommutatives Monoid}{%
    \CommutativeMonoidAlg{A}{\star}{e}}
}

\FormulaAxiomDeltaK[Zugriff auf das Monoid]{%
  \CommutativeMonoidAlg{A}{\star}{e}
  \vdash\MonoidAlg{A}{\star}{e}%
}{CommutativeMonoidBaseMonoidAxiom}{%
  \DeltaRow{Mengen}{A\dsep e}
}

\FormulaAxiomDeltaK[Zugriff auf die Kommutativität eines Monoids]{%
  \CommutativeMonoidAlg{A}{\star}{e}\dsep x,y\in A
  \vdash x\star y=y\star x%
}{CommutativeMonoidCommutativityAxiom}{%
  \DeltaRow{Mengen}{A\dsep e\dsep x\dsep y}
}

\FormulaThmDeltaK[Kommutative Monoide sind kommutative Halbgruppen]{%
  \CommutativeMonoidAlg{A}{\star}{e}
  \vdash\CommutativeSemigroupAlg{A}{\star}%
}{CommutativeMonoidIsCommutativeSemigroup}{%
  \DeltaRow{Mengen}{A\dsep e\dsep x\dsep y}
}
\begin{tabproof}
  \proofstep{1}{\CommutativeMonoidAlg{A}{\star}{e}}{\rA}
  \proofstep{1}{\MonoidAlg{A}{\star}{e}}{%
    \FormulaRefAuto{CommutativeMonoidBaseMonoidAxiom}[ax]{1}}
  \proofstep{1}{\SemigroupAlg{A}{\star}}{%
    \FormulaRefAuto{MonoidBaseSemigroupAxiom}[ax]{2}}
  \proofstep{4}{x\in A}{\rA}
  \proofstep{5}{y\in A}{\rA}
  \proofstep{1,4,5}{x\star y=y\star x}{%
    \FormulaRefAuto{CommutativeMonoidCommutativityAxiom}[ax]{1,4,5}}
  \proofstep{1}{\CommutativeSemigroupAlg{A}{\star}}{%
    \FormulaRefAuto{CommutativeSemigroupDef}[def]{3,6}}
\end{tabproof}

\section{Einheiten, Teilbarkeit, Irreduzibilität und Primelemente}

\begin{remark}
Die in Band~15 eingeführten Halbgruppenbegriffe spezialisieren sich hier auf
die vertrauten Monoidbegriffe.  Ist \(e\) das neutrale Element, so ist
\(\SemigroupUnit{A,\star}{u}\) genau dann erfüllt, wenn
\(\MonoidUnit{A,\star,e}{u}\) gilt.  In einem kommutativen Monoid stimmen
ferner die links-, rechts- und zweiseitige Halbgruppenteilbarkeit mit der
nachstehenden Monoidteilbarkeit überein: Die in Band~15 ausdrücklich
aufgenommene Alternative \(a=b\) wird hier durch den Faktor \(e\) erfasst.
Auch die Echtheitsforderung für ein Halbgruppenprimelement ist dann
äquivalent dazu, dass das Element keine Einheit ist.  Die folgenden
Definitionen sind somit Spezialisierungen, keine konkurrierenden Begriffe.
\end{remark}

\FormulaDefDeltaK[Einheit eines Monoids]{%
  \MonoidUnit{A,\star,e}{u}%
}{MonoidUnitDef}{%
  \DeltaRow{Mengen}{A\dsep e\dsep u\dsep v}
  \DeltaRow{Funktionen}{\star}
  \DeltaPrem{Monoid}{\MonoidAlg{A}{\star}{e}}
  \DeltaRow{Träger}{u\in A}
  \DeltaRow{Inverses Element}{%
    \exists v\in A\,(u\star v=e\land v\star u=e)}
  \DeltaRow{\textbf{Neues Symbol}}{}
  \DeltaPrem{Einheit}{\MonoidUnit{A,\star,e}{u}}
}

\FormulaThmDeltaK[Das neutrale Element ist eine Einheit]{%
  \MonoidAlg{A}{\star}{e}
  \vdash\MonoidUnit{A,\star,e}{e}%
}{MonoidIdentityIsUnit}{%
  \DeltaRow{Mengen}{A\dsep e\dsep v}
  \DeltaRow{Funktionen}{\star}
}
\begin{tabproof}
  \proofstep{1}{\MonoidAlg{A}{\star}{e}}{\rA}
  \proofstep{1}{e\in A}{%
    \FormulaRefAuto{MonoidIdentityCarrierAxiom}[ax]{1}}
  \proofstep{1}{e\star e=e}{%
    \FormulaRefAuto{MonoidLeftIdentityAxiom}[ax]{1,2}}
  \proofstep{1}{%
    \exists v\in A\,(e\star v=e\land v\star e=e)}{%
    \rEI{2,\rAI{3,3}}}
  \proofstep{1}{\MonoidUnit{A,\star,e}{e}}{%
    \FormulaRefAuto{MonoidUnitDef}[def]{1,2,4}}
\end{tabproof}

\paragraph{Beispiel.}
Das neutrale Element muss nicht die einzige Einheit sein. Im additiven Monoid
\((\mathbb Z,+,0_{\mathbb Z})\) ist etwa \(1_{\mathbb Z}\) eine vom
neutralen Element verschiedene Einheit. Ihr inverses Element ist
\(-1_{\mathbb Z}\), denn nach \FormulaRefAuto{IntAddInverse}[thm] gilt
\[
  1_{\mathbb Z}+(-1_{\mathbb Z})=0_{\mathbb Z}
  \qquad\text{und}\qquad
  (-1_{\mathbb Z})+1_{\mathbb Z}=0_{\mathbb Z}.
\]
Dabei ist \(1_{\mathbb Z}\neq0_{\mathbb Z}\). Tatsächlich ist im additiven
Monoid der ganzen Zahlen jede ganze Zahl eine Einheit.

\FormulaDefDeltaK[Teilbarkeit in einem kommutativen Monoid]{%
  \MonoidDivides{A,\star,e}{a}{b}%
}{CommutativeMonoidDividesDef}{%
  \DeltaRow{Mengen}{A\dsep a\dsep b\dsep c\dsep e}
  \DeltaRow{Funktionen}{\star}
  \DeltaPrem{Kommutatives Monoid}{%
    \CommutativeMonoidAlg{A}{\star}{e}}
  \DeltaRow{Elemente}{a,b\in A}
  \DeltaRow{Quotient}{\exists c\in A\,(b=a\star c)}
  \DeltaRow{\textbf{Neues Symbol}}{}
  \DeltaPrem{Teilbarkeit}{\MonoidDivides{A,\star,e}{a}{b}}
}

\FormulaThmDeltaK[Reflexivität der Teilbarkeit]{%
  \CommutativeMonoidAlg{A}{\star}{e}\dsep a\in A
  \vdash\MonoidDivides{A,\star,e}{a}{a}%
}{CommutativeMonoidDividesReflexive}{%
  \DeltaRow{Mengen}{A\dsep a\dsep c\dsep e}
  \DeltaRow{Funktionen}{\star}
}
\begin{tabproof}
  \proofstep{1}{\CommutativeMonoidAlg{A}{\star}{e}}{\rA}
  \proofstep{2}{a\in A}{\rA}
  \proofstep{1}{\MonoidAlg{A}{\star}{e}}{%
    \FormulaRefAuto{CommutativeMonoidBaseMonoidAxiom}[ax]{1}}
  \proofstep{1}{e\in A}{%
    \FormulaRefAuto{MonoidIdentityCarrierAxiom}[ax]{3}}
  \proofstep{1,2}{a\star e=a}{%
    \FormulaRefAuto{MonoidRightIdentityAxiom}[ax]{3,2}}
  \proofstep{1,2}{a=a\star e}{%
    \FormulaRefAuto{a=b\vdash b=a}{5}}
  \proofstep{1,2}{\exists c\in A\,(a=a\star c)}{%
    \rEI{4,6}}
  \proofstep{1,2}{\MonoidDivides{A,\star,e}{a}{a}}{%
    \FormulaRefAuto{CommutativeMonoidDividesDef}[def]{1,2,2,7}}
\end{tabproof}

\FormulaThmDeltaK[Transitivität der Teilbarkeit]{%
  \begin{aligned}[t]
    &\CommutativeMonoidAlg{A}{\star}{e}\dsep a,b,c\in A\dsep
      \MonoidDivides{A,\star,e}{a}{b}\dsep
      \MonoidDivides{A,\star,e}{b}{c}\vdash{}\\[-2pt]
    &\MonoidDivides{A,\star,e}{a}{c}
  \end{aligned}
}{CommutativeMonoidDividesTransitive}{%
  \DeltaRow{Mengen}{A\dsep a\dsep b\dsep c\dsep e\dsep r\dsep s}
  \DeltaRow{Funktionen}{\star}
}
\begin{tabproofsplitwide}
  \proofpartwideindR[Feste Teilbarkeitszeugen verknüpfen]{%
    \begin{aligned}[t]
      &\CommutativeMonoidAlg{A}{\star}{e}\dsep
        a,b,c,r,s\in A\dsep b=a\star r\dsep c=b\star s\vdash{}\\[-2pt]
      &\MonoidDivides{A,\star,e}{a}{c}
    \end{aligned}
  }{%
    \begin{aligned}[t]
      &\CommutativeMonoidAlg{A}{\star}{e}\dsep
        a,b,c,r,s\in A\dsep b=a\star r\dsep c=b\star s\vdash{}\\[-2pt]
      &\MonoidDivides{A,\star,e}{a}{c}
    \end{aligned}
  }[CommutativeMonoidDividesTransitiveFixedWitnesses]
    \proofstepwidestar[1]{\CommutativeMonoidAlg{A}{\star}{e}}{\rA}
    \proofstepwidestar[2]{a\in A}{\rA}
    \proofstepwidestar[3]{b\in A}{\rA}
    \proofstepwidestar[4]{c\in A}{\rA}
    \proofstepwidestar[5]{r\in A}{\rA}
    \proofstepwidestar[6]{s\in A}{\rA}
    \proofstepwidestar[7]{b=a\star r}{\rA}
    \proofstepwidestar[8]{c=b\star s}{\rA}
    \proofstepwidestar[1]{\MonoidAlg{A}{\star}{e}}{%
      \FormulaRefAuto{CommutativeMonoidBaseMonoidAxiom}[ax]{1}}
    \proofstepwidestar[1]{\SemigroupAlg{A}{\star}}{%
      \FormulaRefAuto{MonoidBaseSemigroupAxiom}[ax]{9}}
    \proofstepwidestar[1,5,6]{r\star s\in A}{%
      \FormulaRefAuto{SemigroupClosureAxiom}[ax]{10,5,6}}
    \proofstepwide[8]{c}{=}{b\star s}{\rA}
    \proofstepwide[6,7]{}{=}{(a\star r)\star s}{\rIE{7}}
    \proofstepwide[1,2,5,6]{}{=}{a\star(r\star s)}{%
      \FormulaRefAuto{SemigroupAssociativityAxiom}[ax]{10,2,5,6}}
    \proofstepwidestar[1,2,5,6,7,8]{c=a\star(r\star s)}{%
      \rChain{12,14}}
    \proofstepwidestar[1,2,5,6,7,8]{%
      \exists r\in A\,(c=a\star r)}{\rEI{11,15}}
    \proofstepwidestar[1,2,4,5,6,7,8]{%
      \MonoidDivides{A,\star,e}{a}{c}}{%
      \FormulaRefAuto{CommutativeMonoidDividesDef}[def]{1,2,4,16}}
  \closeproofpartwideindR

  \proofpartwide{Elimination der Teilbarkeitszeugen}
    \proofstepwidestar[1]{\CommutativeMonoidAlg{A}{\star}{e}}{\rA}
    \proofstepwidestar[2]{a\in A}{\rA}
    \proofstepwidestar[3]{b\in A}{\rA}
    \proofstepwidestar[4]{c\in A}{\rA}
    \proofstepwidestar[5]{\MonoidDivides{A,\star,e}{a}{b}}{\rA}
    \proofstepwidestar[6]{\MonoidDivides{A,\star,e}{b}{c}}{\rA}
    \proofstepwidestar[5]{\exists r\in A\,(b=a\star r)}{%
      \FormulaRefAuto{CommutativeMonoidDividesDef}[def]{5}}
    \proofstepwidestar[6]{\exists s\in A\,(c=b\star s)}{%
      \FormulaRefAuto{CommutativeMonoidDividesDef}[def]{6}}
    \proofstepwidestar[9]{r\in A}{\rA}
    \proofstepwidestar[10]{b=a\star r}{\rA}
    \proofstepwidestar[11]{s\in A}{\rA}
    \proofstepwidestar[12]{c=b\star s}{\rA}
    \proofstepwidestar[1,2,3,4,9,10,11,12]{%
      \MonoidDivides{A,\star,e}{a}{c}}{%
      \FormulaRefAuto{CommutativeMonoidDividesTransitiveFixedWitnesses}[thm]{%
        1,2,3,4,9,11,10,12}}
    \proofstepwidestar[1,2,3,4,5,6]{%
      \MonoidDivides{A,\star,e}{a}{c}}{%
      \rEE{7,9,10,\rEE{8,11,12,13}}}
  \closeproofpartwide
\end{tabproofsplitwide}

\FormulaThmDeltaK[Einheiten sind genau die Teiler des neutralen Elements]{%
  \begin{aligned}[t]
    &\CommutativeMonoidAlg{A}{\star}{e}\dsep u\in A\vdash{}\\[-2pt]
    &\MonoidUnit{A,\star,e}{u}
      \leftrightarrow\MonoidDivides{A,\star,e}{u}{e}
  \end{aligned}
}{CommutativeMonoidUnitIffDividesIdentity}{%
  \DeltaRow{Mengen}{A\dsep e\dsep u\dsep v\dsep w}
  \DeltaRow{Funktionen}{\star}
}
\begin{tabproofsplitwide}
  \proofpartwideindR[Eine Einheit teilt das neutrale Element]{%
    \begin{aligned}[t]
      &\CommutativeMonoidAlg{A}{\star}{e}\dsep u\in A\dsep
        \MonoidUnit{A,\star,e}{u}\vdash{}\\[-2pt]
      &\MonoidDivides{A,\star,e}{u}{e}
    \end{aligned}
  }{%
    \begin{aligned}[t]
      &\CommutativeMonoidAlg{A}{\star}{e}\dsep u\in A\dsep
        \MonoidUnit{A,\star,e}{u}\vdash{}\\[-2pt]
      &\MonoidDivides{A,\star,e}{u}{e}
    \end{aligned}
  }[CommutativeMonoidUnitImpliesDividesIdentity]
    \proofstepwidestar[1]{\CommutativeMonoidAlg{A}{\star}{e}}{\rA}
    \proofstepwidestar[2]{u\in A}{\rA}
    \proofstepwidestar[3]{\MonoidUnit{A,\star,e}{u}}{\rA}
    \proofstepwidestar[1]{\MonoidAlg{A}{\star}{e}}{%
      \FormulaRefAuto{CommutativeMonoidBaseMonoidAxiom}[ax]{1}}
    \proofstepwidestar[1]{e\in A}{%
      \FormulaRefAuto{MonoidIdentityCarrierAxiom}[ax]{4}}
    \proofstepwidestar[3]{%
      \exists v\in A\,(u\star v=e\land v\star u=e)}{%
      \FormulaRefAuto{MonoidUnitDef}[def]{3}}
    \proofstepwidestar[7]{v\in A}{\rA}
    \proofstepwidestar[8]{u\star v=e\land v\star u=e}{\rA}
    \proofstepwidestar[8]{u\star v=e}{\rAEa{8}}
    \proofstepwidestar[8]{e=u\star v}{%
      \FormulaRefAuto{a=b\vdash b=a}{9}}
    \proofstepwidestar[7,8]{\exists v\in A\,(e=u\star v)}{%
      \rEI{7,10}}
    \proofstepwidestar[1,2,7,8]{%
      \MonoidDivides{A,\star,e}{u}{e}}{%
      \FormulaRefAuto{CommutativeMonoidDividesDef}[def]{1,2,5,11}}
    \proofstepwidestar[1,2,3]{\MonoidDivides{A,\star,e}{u}{e}}{%
      \rEE{6,7,8,12}}
  \closeproofpartwideindR

  \proofpartwideindR[Ein Teiler des neutralen Elements ist eine Einheit]{%
    \begin{aligned}[t]
      &\CommutativeMonoidAlg{A}{\star}{e}\dsep u\in A\dsep
        \MonoidDivides{A,\star,e}{u}{e}\vdash{}\\[-2pt]
      &\MonoidUnit{A,\star,e}{u}
    \end{aligned}
  }{%
    \begin{aligned}[t]
      &\CommutativeMonoidAlg{A}{\star}{e}\dsep u\in A\dsep
        \MonoidDivides{A,\star,e}{u}{e}\vdash{}\\[-2pt]
      &\MonoidUnit{A,\star,e}{u}
    \end{aligned}
  }[CommutativeMonoidDividesIdentityImpliesUnit]
    \proofstepwidestar[1]{\CommutativeMonoidAlg{A}{\star}{e}}{\rA}
    \proofstepwidestar[2]{u\in A}{\rA}
    \proofstepwidestar[3]{\MonoidDivides{A,\star,e}{u}{e}}{\rA}
    \proofstepwidestar[1]{\MonoidAlg{A}{\star}{e}}{%
      \FormulaRefAuto{CommutativeMonoidBaseMonoidAxiom}[ax]{1}}
    \proofstepwidestar[1]{e\in A}{%
      \FormulaRefAuto{MonoidIdentityCarrierAxiom}[ax]{4}}
    \proofstepwidestar[3]{\exists w\in A\,(e=u\star w)}{%
      \FormulaRefAuto{CommutativeMonoidDividesDef}[def]{3}}
    \proofstepwidestar[7]{w\in A}{\rA}
    \proofstepwidestar[8]{e=u\star w}{\rA}
    \proofstepwidestar[8]{u\star w=e}{%
      \FormulaRefAuto{a=b\vdash b=a}{8}}
    \proofstepwidestar[1,2,7]{w\star u=u\star w}{%
      \FormulaRefAuto{CommutativeMonoidCommutativityAxiom}[ax]{1,7,2}}
    \proofstepwidestar[1,2,7,8]{w\star u=e}{\rChain{10,9}}
    \proofstepwidestar[1,2,7,8]{%
      \exists w\in A\,(u\star w=e\land w\star u=e)}{%
      \rEI{7,\rAI{9,11}}}
    \proofstepwidestar[1,2,7,8]{\MonoidUnit{A,\star,e}{u}}{%
      \FormulaRefAuto{MonoidUnitDef}[def]{4,2,12}}
    \proofstepwidestar[1,2,3]{\MonoidUnit{A,\star,e}{u}}{%
      \rEE{6,7,8,13}}
  \closeproofpartwideindR

  \proofpartwide{Zusammenführung der Richtungen}
    \proofstepwidestar[1]{\CommutativeMonoidAlg{A}{\star}{e}}{\rA}
    \proofstepwidestar[2]{u\in A}{\rA}
    \proofstepwidestar[3]{\MonoidUnit{A,\star,e}{u}}{\rA}
    \proofstepwidestar[1,2,3]{\MonoidDivides{A,\star,e}{u}{e}}{%
      \FormulaRefAuto{CommutativeMonoidUnitImpliesDividesIdentity}[thm]{1,2,3}}
    \proofstepwidestar[1,2]{%
      \MonoidUnit{A,\star,e}{u}\rightarrow
      \MonoidDivides{A,\star,e}{u}{e}}{\rRI{3,4}}
    \proofstepwidestar[6]{\MonoidDivides{A,\star,e}{u}{e}}{\rA}
    \proofstepwidestar[1,2,6]{\MonoidUnit{A,\star,e}{u}}{%
      \FormulaRefAuto{CommutativeMonoidDividesIdentityImpliesUnit}[thm]{1,2,6}}
    \proofstepwidestar[1,2]{%
      \MonoidDivides{A,\star,e}{u}{e}\rightarrow
      \MonoidUnit{A,\star,e}{u}}{\rRI{6,7}}
    \proofstepwidestar[1,2]{%
      \MonoidUnit{A,\star,e}{u}
      \leftrightarrow\MonoidDivides{A,\star,e}{u}{e}}{%
      \rLRI{5,8}}
  \closeproofpartwide
\end{tabproofsplitwide}

\FormulaDefDeltaK[Irreduzibles Element eines kommutativen Monoids]{%
  \MonoidIrreducible{A,\star,e}{p}%
}{CommutativeMonoidIrreducibleDef}{%
  \DeltaRow{Mengen}{A\dsep a\dsep b\dsep e\dsep p}
  \DeltaRow{Funktionen}{\star}
  \DeltaPrem{Kommutatives Monoid}{%
    \CommutativeMonoidAlg{A}{\star}{e}}
  \DeltaRow{Element}{p\in A}
  \DeltaRow{Keine Einheit}{\neg\MonoidUnit{A,\star,e}{p}}
  \DeltaRow{Irreduzibilität}{%
    \begin{aligned}[t]
      &a,b\in A\dsep p=a\star b\vdash{}\\[-2pt]
      &\MonoidUnit{A,\star,e}{a}\lor
       \MonoidUnit{A,\star,e}{b}
    \end{aligned}}
  \DeltaRow{\textbf{Neues Symbol}}{}
  \DeltaPrem{Irreduzibles Element}{%
    \MonoidIrreducible{A,\star,e}{p}}
}

\FormulaDefDeltaK[Primelement eines kommutativen Monoids]{%
  \MonoidPrime{A,\star,e}{p}%
}{CommutativeMonoidPrimeDef}{%
  \DeltaRow{Mengen}{A\dsep a\dsep b\dsep e\dsep p}
  \DeltaRow{Funktionen}{\star}
  \DeltaPrem{Kommutatives Monoid}{%
    \CommutativeMonoidAlg{A}{\star}{e}}
  \DeltaRow{Element}{p\in A}
  \DeltaRow{Keine Einheit}{\neg\MonoidUnit{A,\star,e}{p}}
  \DeltaRow{Primeigenschaft}{%
    \begin{aligned}[t]
      &a,b\in A\dsep
        \MonoidDivides{A,\star,e}{p}{a\star b}\vdash{}\\[-2pt]
      &\MonoidDivides{A,\star,e}{p}{a}\lor
       \MonoidDivides{A,\star,e}{p}{b}
    \end{aligned}}
  \DeltaRow{\textbf{Neues Symbol}}{}
  \DeltaPrem{Primelement}{\MonoidPrime{A,\star,e}{p}}
}

\FormulaThmDeltaK[Ein geteilter Faktor erzwingt eine Einheit]{%
  \begin{aligned}[t]
    &\CommutativeMonoidAlg{A}{\star}{e}\dsep
      \mathsf{K\ddot{u}rz}\!\left(A,\star\right)\dsep
      p,a,b\in A\dsep p=a\star b\dsep{}\\[-2pt]
    &\MonoidDivides{A,\star,e}{p}{a}
      \vdash\MonoidUnit{A,\star,e}{b}
  \end{aligned}
}{CommutativeCancellativeMonoidFactorDivisorUnit}{%
  \DeltaRow{Mengen}{A\dsep a\dsep b\dsep c\dsep e\dsep p}
  \DeltaRow{Funktionen}{\star}
}
\begin{tabproofsplitwide}
  \proofpartwideindR[Herstellung der Kürzungsgleichung]{%
    \begin{aligned}[t]
      &\CommutativeMonoidAlg{A}{\star}{e}\dsep
        p,a,b,c\in A\dsep p=a\star b\dsep a=p\star c\vdash{}\\[-2pt]
      &p\star e=p\star(c\star b)
    \end{aligned}
  }{%
    \begin{aligned}[t]
      &\CommutativeMonoidAlg{A}{\star}{e}\dsep
        p,a,b,c\in A\dsep p=a\star b\dsep a=p\star c\vdash{}\\[-2pt]
      &p\star e=p\star(c\star b)
    \end{aligned}
  }[CommutativeMonoidFactorDivisorCancellationEquation]
    \proofstepwidestar[1]{\CommutativeMonoidAlg{A}{\star}{e}}{\rA}
    \proofstepwidestar[2]{p\in A}{\rA}
    \proofstepwidestar[3]{a\in A}{\rA}
    \proofstepwidestar[4]{b\in A}{\rA}
    \proofstepwidestar[5]{c\in A}{\rA}
    \proofstepwidestar[6]{p=a\star b}{\rA}
    \proofstepwidestar[7]{a=p\star c}{\rA}
    \proofstepwidestar[1]{\MonoidAlg{A}{\star}{e}}{%
      \FormulaRefAuto{CommutativeMonoidBaseMonoidAxiom}[ax]{1}}
    \proofstepwidestar[1]{\SemigroupAlg{A}{\star}}{%
      \FormulaRefAuto{MonoidBaseSemigroupAxiom}[ax]{8}}
    \proofstepwide[1,2]{p\star e}{=}{p}{%
      \FormulaRefAuto{MonoidRightIdentityAxiom}[ax]{8,2}}
    \proofstepwide[6]{}{=}{a\star b}{\rA}
    \proofstepwide[4,7]{}{=}{(p\star c)\star b}{\rIE{7}}
    \proofstepwide[1,2,4,5]{}{=}{p\star(c\star b)}{%
      \FormulaRefAuto{SemigroupAssociativityAxiom}[ax]{9,2,5,4}}
    \proofstepwidestar[1,2,4,5,6,7]{%
      p\star e=p\star(c\star b)}{\rChain{10,13}}
  \closeproofpartwideindR

  \proofpartwideindR[Kürzung und Gewinnung des Inversen]{%
    \begin{aligned}[t]
      &\CommutativeMonoidAlg{A}{\star}{e}\dsep
        \mathsf{K\ddot{u}rz}\!\left(A,\star\right)\dsep
        p,b,c\in A\dsep{}\\[-2pt]
      &p\star e=p\star(c\star b)
        \vdash\MonoidUnit{A,\star,e}{b}
    \end{aligned}
  }{%
    \begin{aligned}[t]
      &\CommutativeMonoidAlg{A}{\star}{e}\dsep
        \mathsf{K\ddot{u}rz}\!\left(A,\star\right)\dsep
        p,b,c\in A\dsep{}\\[-2pt]
      &p\star e=p\star(c\star b)
        \vdash\MonoidUnit{A,\star,e}{b}
    \end{aligned}
  }[CommutativeCancellativeMonoidCancellationEquationImpliesUnit]
    \proofstepwidestar[1]{\CommutativeMonoidAlg{A}{\star}{e}}{\rA}
    \proofstepwidestar[2]{\mathsf{K\ddot{u}rz}\!\left(A,\star\right)}{\rA}
    \proofstepwidestar[3]{p\in A}{\rA}
    \proofstepwidestar[4]{b\in A}{\rA}
    \proofstepwidestar[5]{c\in A}{\rA}
    \proofstepwidestar[6]{p\star e=p\star(c\star b)}{\rA}
    \proofstepwidestar[1]{\MonoidAlg{A}{\star}{e}}{%
      \FormulaRefAuto{CommutativeMonoidBaseMonoidAxiom}[ax]{1}}
    \proofstepwidestar[1]{\SemigroupAlg{A}{\star}}{%
      \FormulaRefAuto{MonoidBaseSemigroupAxiom}[ax]{7}}
    \proofstepwidestar[1]{e\in A}{%
      \FormulaRefAuto{MonoidIdentityCarrierAxiom}[ax]{7}}
    \proofstepwidestar[2]{\mathsf{LK\ddot{u}rz}\!\left(A,\star\right)}{%
      \FormulaRefAuto{CancellativeLeftAxiom}[ax]{2}}
    \proofstepwidestar[1,4,5]{c\star b\in A}{%
      \FormulaRefAuto{SemigroupClosureAxiom}[ax]{8,5,4}}
    \proofstepwidestar[1,2,3,4,5,6]{e=c\star b}{%
      \FormulaRefAuto{LeftCancellationAxiom}[ax]{10,3,9,11,6}}
    \proofstepwidestar[1,2,3,4,5,6]{c\star b=e}{%
      \FormulaRefAuto{a=b\vdash b=a}{12}}
    \proofstepwide[1,4,5]{b\star c}{=}{c\star b}{%
      \FormulaRefAuto{CommutativeMonoidCommutativityAxiom}[ax]{1,4,5}}
    \proofstepwidestar[1,2,3,4,5,6]{b\star c=e}{%
      \rChain{14,13}}
    \proofstepwidestar[1,2,3,4,5,6]{%
      \exists c\in A\,(b\star c=e\land c\star b=e)}{%
      \rEI{5,\rAI{15,13}}}
    \proofstepwidestar[1,2,3,4,5,6]{\MonoidUnit{A,\star,e}{b}}{%
      \FormulaRefAuto{MonoidUnitDef}[def]{7,4,16}}
  \closeproofpartwideindR

  \proofpartwide{Elimination des Teilbarkeitszeugen}
    \proofstepwidestar[1]{\CommutativeMonoidAlg{A}{\star}{e}}{\rA}
    \proofstepwidestar[2]{\mathsf{K\ddot{u}rz}\!\left(A,\star\right)}{\rA}
    \proofstepwidestar[3]{p\in A}{\rA}
    \proofstepwidestar[4]{a\in A}{\rA}
    \proofstepwidestar[5]{b\in A}{\rA}
    \proofstepwidestar[6]{p=a\star b}{\rA}
    \proofstepwidestar[7]{\MonoidDivides{A,\star,e}{p}{a}}{\rA}
    \proofstepwidestar[7]{\exists c\in A\,(a=p\star c)}{%
      \FormulaRefAuto{CommutativeMonoidDividesDef}[def]{7}}
    \proofstepwidestar[9]{c\in A}{\rA}
    \proofstepwidestar[10]{a=p\star c}{\rA}
    \proofstepwidestar[1,3,4,5,6,9,10]{p\star e=p\star(c\star b)}{%
      \FormulaRefAuto{CommutativeMonoidFactorDivisorCancellationEquation}[thm]{%
        1,3,4,5,9,6,10}}
    \proofstepwidestar[1,2,3,4,5,6,9,10]{\MonoidUnit{A,\star,e}{b}}{%
      \FormulaRefAuto{CommutativeCancellativeMonoidCancellationEquationImpliesUnit}[thm]{%
        1,2,3,5,9,11}}
    \proofstepwidestar[1,2,3,4,5,6,7]{\MonoidUnit{A,\star,e}{b}}{%
      \rEE{8,9,10,12}}
  \closeproofpartwide
\end{tabproofsplitwide}

\FormulaThmDeltaK[Primelemente in kommutativen kürzbaren Monoiden sind irreduzibel]{%
  \begin{aligned}[t]
    &\CommutativeMonoidAlg{A}{\star}{e}\dsep
      \mathsf{K\ddot{u}rz}\!\left(A,\star\right)\dsep
      \MonoidPrime{A,\star,e}{p}\vdash{}\\[-2pt]
    &\MonoidIrreducible{A,\star,e}{p}
  \end{aligned}
}{CommutativeCancellativeMonoidPrimeIsIrreducible}{%
  \DeltaRow{Mengen}{A\dsep a\dsep b\dsep e\dsep p}
  \DeltaRow{Funktionen}{\star}
}
\begin{tabproofsplitwide}
  \proofpartwideindR[Teilerdisjunktion der Faktorisierung]{%
    \begin{aligned}[t]
      &\CommutativeMonoidAlg{A}{\star}{e}\dsep
        \MonoidPrime{A,\star,e}{p}\dsep{}\\[-2pt]
      &a,b\in A\dsep p=a\star b\vdash{}\\[-2pt]
      &\MonoidDivides{A,\star,e}{p}{a}\lor
       \MonoidDivides{A,\star,e}{p}{b}
    \end{aligned}
  }{%
    \begin{aligned}[t]
      &\CommutativeMonoidAlg{A}{\star}{e}\dsep
        \MonoidPrime{A,\star,e}{p}\dsep{}\\[-2pt]
      &a,b\in A\dsep p=a\star b\vdash{}\\[-2pt]
      &\MonoidDivides{A,\star,e}{p}{a}\lor
       \MonoidDivides{A,\star,e}{p}{b}
    \end{aligned}
  }[CommutativeMonoidPrimeFactorDivisorDisjunction]
    \proofstepwidestar[1]{\CommutativeMonoidAlg{A}{\star}{e}}{\rA}
    \proofstepwidestar[2]{\MonoidPrime{A,\star,e}{p}}{\rA}
    \proofstepwidestar[3]{a\in A}{\rA}
    \proofstepwidestar[4]{b\in A}{\rA}
    \proofstepwidestar[5]{p=a\star b}{\rA}
    \proofstepwidestar[2]{p\in A}{%
      \FormulaRefAuto{CommutativeMonoidPrimeDef}[def]{2}}
    \proofstepwidestar[1,2]{\MonoidDivides{A,\star,e}{p}{p}}{%
      \FormulaRefAuto{CommutativeMonoidDividesReflexive}[thm]{1,6}}
    \proofstepwidestar[1,2,5]{%
      \MonoidDivides{A,\star,e}{p}{a\star b}}{\rIE{5}}
    \proofstepwidestar[1,2,3,4,5]{%
      \MonoidDivides{A,\star,e}{p}{a}\lor
      \MonoidDivides{A,\star,e}{p}{b}}{%
      \FormulaRefAuto{CommutativeMonoidPrimeDef}[def]{2,3,4,8}}
  \closeproofpartwideindR

  \proofpartwideindR[Von der Teilerdisjunktion zur Einheitendisjunktion]{%
    \begin{aligned}[t]
      &\CommutativeMonoidAlg{A}{\star}{e}\dsep
        \mathsf{K\ddot{u}rz}\!\left(A,\star\right)\dsep
        \MonoidPrime{A,\star,e}{p}\dsep a,b\in A\dsep{}\\[-2pt]
      &p=a\star b\dsep
        \bigl(\MonoidDivides{A,\star,e}{p}{a}\lor
        \MonoidDivides{A,\star,e}{p}{b}\bigr)\vdash{}\\[-2pt]
      &\MonoidUnit{A,\star,e}{a}\lor\MonoidUnit{A,\star,e}{b}
    \end{aligned}
  }{%
    \begin{aligned}[t]
      &\CommutativeMonoidAlg{A}{\star}{e}\dsep
        \mathsf{K\ddot{u}rz}\!\left(A,\star\right)\dsep
        \MonoidPrime{A,\star,e}{p}\dsep a,b\in A\dsep{}\\[-2pt]
      &p=a\star b\dsep
        \bigl(\MonoidDivides{A,\star,e}{p}{a}\lor
        \MonoidDivides{A,\star,e}{p}{b}\bigr)\vdash{}\\[-2pt]
      &\MonoidUnit{A,\star,e}{a}\lor\MonoidUnit{A,\star,e}{b}
    \end{aligned}
  }[CommutativeCancellativeMonoidPrimeFactorUnitDisjunction]
    \proofstepwidestar[1]{\CommutativeMonoidAlg{A}{\star}{e}}{\rA}
    \proofstepwidestar[2]{\mathsf{K\ddot{u}rz}\!\left(A,\star\right)}{\rA}
    \proofstepwidestar[3]{\MonoidPrime{A,\star,e}{p}}{\rA}
    \proofstepwidestar[4]{a\in A}{\rA}
    \proofstepwidestar[5]{b\in A}{\rA}
    \proofstepwidestar[6]{p=a\star b}{\rA}
    \proofstepwidestar[7]{%
      \MonoidDivides{A,\star,e}{p}{a}\lor
      \MonoidDivides{A,\star,e}{p}{b}}{\rA}
    \proofstepwidestar[3]{p\in A}{%
      \FormulaRefAuto{CommutativeMonoidPrimeDef}[def]{3}}
    \proofstepwidestar[9]{\MonoidDivides{A,\star,e}{p}{a}}{\rA}
    \proofstepwidestar[1,2,3,4,5,6,9]{\MonoidUnit{A,\star,e}{b}}{%
      \FormulaRefAuto{CommutativeCancellativeMonoidFactorDivisorUnit}[thm]{%
        1,2,8,4,5,6,9}}
    \proofstepwidestar[1,2,3,4,5,6,9]{%
      \MonoidUnit{A,\star,e}{a}\lor\MonoidUnit{A,\star,e}{b}}{%
      \rOIb{10}}
    \proofstepwidestar[12]{\MonoidDivides{A,\star,e}{p}{b}}{\rA}
    \proofstepwidestar[1,4,5]{a\star b=b\star a}{%
      \FormulaRefAuto{CommutativeMonoidCommutativityAxiom}[ax]{1,4,5}}
    \proofstepwidestar[1,4,5,6]{p=b\star a}{\rChain{6,13}}
    \proofstepwidestar[1,2,3,4,5,6,12]{\MonoidUnit{A,\star,e}{a}}{%
      \FormulaRefAuto{CommutativeCancellativeMonoidFactorDivisorUnit}[thm]{%
        1,2,8,5,4,14,12}}
    \proofstepwidestar[1,2,3,4,5,6,12]{%
      \MonoidUnit{A,\star,e}{a}\lor\MonoidUnit{A,\star,e}{b}}{%
      \rOIa{15}}
    \proofstepwidestar[1,2,3,4,5,6,7]{%
      \MonoidUnit{A,\star,e}{a}\lor\MonoidUnit{A,\star,e}{b}}{%
      \rOE{7,9,11,12,16}}
  \closeproofpartwideindR

  \proofpartwide{Nachweis der Irreduzibilität}
    \proofstepwidestar[1]{\CommutativeMonoidAlg{A}{\star}{e}}{\rA}
    \proofstepwidestar[2]{\mathsf{K\ddot{u}rz}\!\left(A,\star\right)}{\rA}
    \proofstepwidestar[3]{\MonoidPrime{A,\star,e}{p}}{\rA}
    \proofstepwidestar[3]{p\in A}{%
      \FormulaRefAuto{CommutativeMonoidPrimeDef}[def]{3}}
    \proofstepwidestar[3]{\neg\MonoidUnit{A,\star,e}{p}}{%
      \FormulaRefAuto{CommutativeMonoidPrimeDef}[def]{3}}
    \proofstepwidestar[6]{a\in A}{\rA}
    \proofstepwidestar[7]{b\in A}{\rA}
    \proofstepwidestar[8]{p=a\star b}{\rA}
    \proofstepwidestar[1,3,6,7,8]{%
      \MonoidDivides{A,\star,e}{p}{a}\lor
      \MonoidDivides{A,\star,e}{p}{b}}{%
      \FormulaRefAuto{CommutativeMonoidPrimeFactorDivisorDisjunction}[thm]{%
        1,3,6,7,8}}
    \proofstepwidestar[1,2,3,6,7,8]{%
      \MonoidUnit{A,\star,e}{a}\lor\MonoidUnit{A,\star,e}{b}}{%
      \FormulaRefAuto{CommutativeCancellativeMonoidPrimeFactorUnitDisjunction}[thm]{%
        1,2,3,6,7,8,9}}
    \proofstepwidestar[1,2,3]{\MonoidIrreducible{A,\star,e}{p}}{%
      \FormulaRefAuto{CommutativeMonoidIrreducibleDef}[def]{%
        1,4,5,6,7,8,10}}
  \closeproofpartwide
\end{tabproofsplitwide}

\section{Morphismen von Monoiden}

\FormulaDefDeltaK[Begriff des Monoidhomomorphismus]{%
  \MonoidHom{f}{A,\star,e}{B,\diamond,u}%
}{MonoidHomDef}{%
  \DeltaRow{Mengen}{A\dsep B\dsep e\dsep u}
  \DeltaRow{Funktionen}{f\dsep\star\dsep\diamond}
  \DeltaRow{\textbf{Axiome}}{}
  \DeltaPrem{Quellmonoid}{\MonoidAlg{A}{\star}{e}}
  \DeltaPrem{Zielmonoid}{\MonoidAlg{B}{\diamond}{u}}
  \DeltaPrem{Halbgruppenhomomorphismus}{%
    \SemigroupHom{f}{A,\star}{B,\diamond}}
  \DeltaRow{Erhaltung des neutralen Elements}{f(e)=u}
  \DeltaRow{\textbf{Neues Symbol}}{}
  \DeltaPrem{Monoidhomomorphismus}{%
    \MonoidHom{f}{A,\star,e}{B,\diamond,u}}
}

\FormulaAxiomDeltaK[Operationserhaltung eines Monoidhomomorphismus]{%
  \MonoidHom{f}{A,\star,e}{B,\diamond,u}\dsep x,y\in A
  \vdash f(x\star y)=f(x)\diamond f(y)%
}{MonoidHomOperationAxiom}{%
  \DeltaRow{Mengen}{A\dsep B\dsep e\dsep u\dsep x\dsep y}
  \DeltaRow{Funktionen}{f\dsep\star\dsep\diamond}
}

\FormulaAxiomDeltaK[Erhaltung des neutralen Elements]{%
  \MonoidHom{f}{A,\star,e}{B,\diamond,u}
  \vdash f(e)=u%
}{MonoidHomIdentityAxiom}{%
  \DeltaRow{Mengen}{A\dsep B\dsep e\dsep u}
  \DeltaRow{Funktionen}{f\dsep\star\dsep\diamond}
}

\FormulaDefDeltaK[Begriff des Monoidisomorphismus]{%
  \MonoidIso{f}{A,\star,e}{B,\diamond,u}%
}{MonoidIsoDef}{%
  \DeltaRow{Mengen}{A\dsep B\dsep e\dsep u}
  \DeltaRow{Funktionen}{f\dsep\star\dsep\diamond}
  \DeltaRow{\textbf{Axiome}}{}
  \DeltaPrem{Monoidhomomorphismus}{%
    \MonoidHom{f}{A,\star,e}{B,\diamond,u}}
  \DeltaPrem{Bijektivität}{f\colon A\bij B}
  \DeltaRow{\textbf{Neues Symbol}}{}
  \DeltaPrem{Monoidisomorphismus}{%
    \MonoidIso{f}{A,\star,e}{B,\diamond,u}}
}

\subsection{Grundlegende Eigenschaften}

\FormulaThmDeltaK[Komposition von Monoidhomomorphismen]{%
  \begin{aligned}[t]
    &\MonoidHom{f}{A,\star,e}{B,\diamond,u}\dsep
      \MonoidHom{g}{B,\diamond,u}{C,\bullet,v}\\
    &\vdash\MonoidHom{g\circ f}{A,\star,e}{C,\bullet,v}
  \end{aligned}
}{MonoidHomComposition}{%
  \DeltaRow{Mengen}{A\dsep B\dsep C\dsep e\dsep u\dsep v}
  \DeltaRow{Funktionen}{f\dsep g\dsep\star\dsep\diamond\dsep\bullet}
}
\begin{tabproofwide}
  \proofstepwidestar[1]{\MonoidHom{f}{A,\star,e}{B,\diamond,u}}{\rA}
  \proofstepwidestar[2]{\MonoidHom{g}{B,\diamond,u}{C,\bullet,v}}{\rA}
  \proofstepwidestar[1]{\SemigroupHom{f}{A,\star}{B,\diamond}}{%
    \FormulaRefAuto{MonoidHomDef}[def]{1}}
  \proofstepwidestar[2]{\SemigroupHom{g}{B,\diamond}{C,\bullet}}{%
    \FormulaRefAuto{MonoidHomDef}[def]{2}}
  \proofstepwidestar[1,2]{%
    \SemigroupHom{g\circ f}{A,\star}{C,\bullet}}{%
    \FormulaRefAuto{SemigroupHomComposition}[thm]{3,4}}
  \proofstepwidestar[1]{f(e)=u}{%
    \FormulaRefAuto{MonoidHomIdentityAxiom}[ax]{1}}
  \proofstepwidestar[2]{g(u)=v}{%
    \FormulaRefAuto{MonoidHomIdentityAxiom}[ax]{2}}
  \proofstepwidestar[1]{\MonoidAlg{A}{\star}{e}}{%
    \FormulaRefAuto{MonoidHomDef}[def]{1}}
  \proofstepwidestar[1]{e\in A}{%
    \FormulaRefAuto{MonoidIdentityCarrierAxiom}[ax]{8}}
  \proofstepwidestar[1]{f\colon A\to B}{%
    \FormulaRefAuto{SemigroupHomFunctionAxiom}[ax]{3}}
  \proofstepwidestar[2]{g\colon B\to C}{%
    \FormulaRefAuto{SemigroupHomFunctionAxiom}[ax]{4}}
  \proofstepwide[1,2]{(g\circ f)(e)}{=}{g(f(e))}{%
    \FormulaRefAuto{x\in A \vdash (G\circ F)(x)=G(F(x))}{%
      10,11,9}}
  \proofstepwide[1,2]{}{=}{g(u)}{\rIE{6}}
  \proofstepwide[1,2]{}{=}{v}{\rIE{7}}
  \proofstepwidestar[1,2]{(g\circ f)(e)=v}{\rChain{12,14}}
  \proofstepwidestar[2]{\MonoidAlg{C}{\bullet}{v}}{%
    \FormulaRefAuto{MonoidHomDef}[def]{2}}
  \proofstepwidestar[1,2]{%
    \MonoidHom{g\circ f}{A,\star,e}{C,\bullet,v}}{%
    \FormulaRefAuto{MonoidHomDef}[def]{8,16,5,15}}
\end{tabproofwide}

\subsection{Kommutative Homomorphismen und Isomorphismen}

\FormulaDefDeltaK[Homomorphismen kommutativer Monoide]{%
  \begin{aligned}[t]
    &\CommutativeMonoidHom{f}{A,\star,e}{B,\diamond,u}\\[-2pt]
    &\quad\coloneqq
      \CommutativeMonoidAlg{A}{\star}{e}\land
      \CommutativeMonoidAlg{B}{\diamond}{u}\\[-2pt]
    &\qquad\land\MonoidHom{f}{A,\star,e}{B,\diamond,u}
  \end{aligned}
}{CommutativeMonoidHomDef}{%
  \DeltaRow{Mengen}{A\dsep B\dsep e\dsep u}
  \DeltaRow{Funktionen}{f\dsep\star\dsep\diamond}
  \DeltaRow{Operationserhaltung}{f(x\star y)=f(x)\diamond f(y)}
  \DeltaRow{Erhaltung des Neutralen}{f(e)=u}
  \DeltaRow{\textbf{Neues Symbol}}{}
  \DeltaPrem{Symbol}{\mathsf{KMonHom}}
}

\FormulaAxiomDeltaK[Quellstruktur eines Homomorphismus kommutativer Monoide]{%
  \CommutativeMonoidHom{f}{A,\star,e}{B,\diamond,u}
  \vdash\CommutativeMonoidAlg{A}{\star}{e}%
}{CommutativeMonoidHomSourceAxiom}{%
  \DeltaRow{Mengen}{A\dsep B\dsep e\dsep u}
  \DeltaRow{Funktionen}{f\dsep\star\dsep\diamond}
}

\FormulaAxiomDeltaK[Zielstruktur eines Homomorphismus kommutativer Monoide]{%
  \CommutativeMonoidHom{f}{A,\star,e}{B,\diamond,u}
  \vdash\CommutativeMonoidAlg{B}{\diamond}{u}%
}{CommutativeMonoidHomTargetAxiom}{%
  \DeltaRow{Mengen}{A\dsep B\dsep e\dsep u}
  \DeltaRow{Funktionen}{f\dsep\star\dsep\diamond}
}

\FormulaAxiomDeltaK[Monoidhomomorphismus eines kommutativen Homomorphismus]{%
  \CommutativeMonoidHom{f}{A,\star,e}{B,\diamond,u}
  \vdash\MonoidHom{f}{A,\star,e}{B,\diamond,u}%
}{CommutativeMonoidHomBaseHomAxiom}{%
  \DeltaRow{Mengen}{A\dsep B\dsep e\dsep u}
  \DeltaRow{Funktionen}{f\dsep\star\dsep\diamond}
}

\FormulaDefDeltaK[Isomorphismen kommutativer Monoide]{%
  \begin{aligned}[t]
    &\CommutativeMonoidIso{f}{A,\star,e}{B,\diamond,u}\\[-2pt]
    &\quad\coloneqq
      \CommutativeMonoidHom{f}{A,\star,e}{B,\diamond,u}
      \land f\colon A\bij B
  \end{aligned}
}{CommutativeMonoidIsoDef}{%
  \DeltaRow{Mengen}{A\dsep B\dsep e\dsep u}
  \DeltaRow{Funktionen}{f\dsep\star\dsep\diamond}
  \DeltaRow{\textbf{Neues Symbol}}{}
  \DeltaPrem{Isomorphismus kommutativer Monoide}{\mathsf{KMonIso}}
}

\FormulaAxiomDeltaK[Homomorphismus eines Isomorphismus kommutativer Monoide]{%
  \CommutativeMonoidIso{f}{A,\star,e}{B,\diamond,u}
  \vdash\CommutativeMonoidHom{f}{A,\star,e}{B,\diamond,u}%
}{CommutativeMonoidIsoHomAxiom}{%
  \DeltaRow{Mengen}{A\dsep B\dsep e\dsep u}
  \DeltaRow{Funktionen}{f\dsep\star\dsep\diamond}
}

\FormulaAxiomDeltaK[Bijektivität eines Isomorphismus kommutativer Monoide]{%
  \CommutativeMonoidIso{f}{A,\star,e}{B,\diamond,u}
  \vdash f\colon A\bij B%
}{CommutativeMonoidIsoBijectiveAxiom}{%
  \DeltaRow{Mengen}{A\dsep B\dsep e\dsep u}
  \DeltaRow{Funktionen}{f\dsep\star\dsep\diamond}
}

\end{document}
